\chapter{证据、来源与时间}

\section{来源对象}

Nous Wave 以可追溯来源作为认知形成的基础。来源层至少包含原始材料、观察事件、派生表示和证据引用。

\begin{description}
\item[原始材料（Artifact）] 稳定的内容实体，可以是消息载荷、文件、网页快照、图像、音频、工具输出或其他材料。
\item[观察事件（ObservationOccurrence）] 主体在特定时间、渠道和上下文中实际接触某份材料或事件的一次记录。同一材料可以对应多次观察事件。
\item[派生表示（DerivedRepresentation）] 由解析器、光学字符识别、模型或其他转换生成的版本化表示，保存生产者、版本、输入来源和区域映射。
\item[证据引用（EvidenceRef）] 指向观察事件、来源区域、派生表示或派生区域的稳定引用。
\end{description}

原始材料的内容去重不得消除观察事件身份。主体两次看到相同内容时，可以共享材料实体，但必须保留两次实际接触。

\section{来源谱系}

所有派生表示保存完整来源链。多个摘要、转写或模型派生如果来自同一原始来源，只构成同一来源祖先的不同表示。系统计算独立佐证时以来源祖先和形成过程为依据，不按文本数量、向量点数量或派生表示数量计数。

生产者签名记录生成派生表示或提案所需的实现身份，包括模型、版本、工作流和关键配置。生产者签名用于重建、失效判断和审计，不成为认知真值本身。

\section{来源提及与认知指向}

来源提及（EntityMention）记录材料某一区域明确出现的实体。认知指向（CognitionAboutness）记录某项认知在语义上涉及的实体、角色或关系。二者分别维护。

认知指向可以由多个来源共同支持，也可以在来源中没有出现规范实体名称时通过其他证据得到。来源提及不能因为某个认知指向该实体而被反向伪造。

\section{时间语义}

系统分别记录以下时间：

\begin{description}
\item[发生时间] 外部事件实际发生的时间或时间范围。
\item[观察时间] 主体实际接触事件或材料的时间。
\item[有效时间] 某项主张在所描述世界中成立的时间范围。
\item[形成时间] 主体形成某项认知的时间。
\item[记录时间] 系统持久化相应记录的时间。
\item[权威提交序号] 同一主体范围内用于确定持久提交顺序的单调序号。
\end{description}

未知时间保持未知。时间范围允许开放端点。读取层可以派生最早、最晚或聚合时间视图，但派生值不写回为新的来源事实。
